%!TEX program = uplatex
\RequirePackage{plautopatch}
\documentclass[uplatex,dvipdfmx]{jlreq}
\usepackage{amsmath,amssymb}

\pagestyle{empty}
\begin{document}

\noindent
{\large\bfseries Log Scheme 理論入門}
\hfill
{\large\bfseries 諏訪紀幸（中央大学・理工学部）}

\bigskip

加藤和也、Fontaine、Illusie によって十年程前導入された
「logarithmic structure を持つ scheme」すなわち log scheme の理論は、
toroidal geometry のひとつの一般化を与えているが、
中山能力、梶原健、中島幸喜、加藤文元、志甫淳といった国内の若手を中心に、
近年目覚しく研究が進んでおり、広く注目を集めている。

本講演では、toroidal geometry との関連に主眼をおき、log scheme 理論の概説を試みる。
大学院生や専門以外の方にも、できるだけ具体的なイメージを持ってもらえるように注意しながら、
\begin{enumerate}
  \setlength{\itemsep}{2pt}
  \item[(1)] log scheme とは何か。それが toroidal geometry の一般化と見なせるのは何故か。
  \item[(2)] どのような応用があるか。
\end{enumerate}
の二点を解説することが目標である。

要点をもう少し詳しく述べよう。
toroidal geometry の理論では扇（fan）が活躍した。
一方、log scheme の理論では半群が活躍する。
scheme の構造層が乗法に関してなす半群に $\mathbf{N}$ のような性質の良い半群を
上手に絡ませて、scheme の morphism が smooth あるいは étale でなくても、
あたかも smooth あるいは étale であるかのような取り扱いが可能になる。
誤解を恐れずに言えば、log scheme の世界では通常２重点を smooth な点と見なすことが出来る。
具体的な応用としては、例えば Tate curve や Mumford curve のような
semi-stable reduction をもつ局所体の上の代数多様体を考える際に、
log scheme の理論は良い枠組みを提供している。

これらの点について、個人的な思い入れも交えてざっくばらんに語ってみたいと考えている。

\end{document}